\textbf{Phần 1. Trắc nghiệm lựa chọn}
\setcounterpageref{ex}{0}

\begin{ex} %[1D8V30]
Tính kết quả của phép tính: $23 + 45$.
\choice
{$67$}
{\True $68$}
{$69$}
{$66$}
\loigiai{
Hàng đơn vị: $3 + 5 = 8$. Hàng chục: $2 + 4 = 6$. Vậy $23 + 45 = 68$.
\textbf{Đáp số:} $68$.\par
\textbf{Đáp án của câu hỏi là B}
}
\end{ex}

\begin{ex} %[1D8V30]
Tính kết quả của phép tính: $34 + 52$.
\choice
{\True $86$}
{$85$}
{$87$}
{$84$}
\loigiai{
Hàng đơn vị: $4 + 2 = 6$. Hàng chục: $3 + 5 = 8$. Vậy $34 + 52 = 86$.
\textbf{Đáp số:} $86$.\par
\textbf{Đáp án của câu hỏi là A}
}
\end{ex}

\begin{ex} %[1D8V30]
Tính kết quả của phép tính: $16 + 71$.
\choice
{$86$}
{$88$}
{\True $87$}
{$85$}
\loigiai{
Hàng đơn vị: $6 + 1 = 7$. Hàng chục: $1 + 7 = 8$. Vậy $16 + 71 = 87$.
\textbf{Đáp số:} $87$.\par
\textbf{Đáp án của câu hỏi là C}
}
\end{ex}

\begin{ex} %[1D8V30]
Điền số thích hợp: $25 + \dots = 69$.
\choice
{$43$}
{$45$}
{$42$}
{\True $44$}
\loigiai{
Số cần điền là: $69 - 25 = 44$.
\textbf{Đáp số:} $44$.\par
\textbf{Đáp án của câu hỏi là D}
}
\end{ex}

\begin{ex} %[1D8V30]
Điền số thích hợp: $\dots + 31 = 85$.
\choice
{$53$}
{\True $54$}
{$55$}
{$52$}
\loigiai{
Số cần điền là: $85 - 31 = 54$.
\textbf{Đáp số:} $54$.\par
\textbf{Đáp án của câu hỏi là B}
}
\end{ex}

\begin{ex} %[1D8V30]
Lớp 1A có 32 học sinh, lớp 1B có 35 học sinh. Hỏi cả hai lớp có tất cả bao nhiêu học sinh?
\choice
{\True $67$}
{$66$}
{$68$}
{$65$}
\loigiai{
Tổng số học sinh cả hai lớp là: $32 + 35 = 67$ (học sinh).
\textbf{Đáp số:} $67$.\par
\textbf{Đáp án của câu hỏi là A}
}
\end{ex}

\begin{ex} %[1D8V30]
Một người có 40 quả cam và 35 quả quýt. Hỏi người đó có tất cả bao nhiêu quả cả hai loại?
\choice
{$74$}
{$76$}
{\True $75$}
{$73$}
\loigiai{
Tổng số quả là: $40 + 35 = 75$ (quả).
\textbf{Đáp số:} $75$.\par
\textbf{Đáp án của câu hỏi là C}
}
\end{ex}

\begin{ex} %[1D8V30]
Bạn Hùng có 24 viên bi, bạn Dũng có 33 viên bi. Hỏi cả hai bạn có bao nhiêu viên bi?
\choice
{$56$}
{$58$}
{$55$}
{\True $57$}
\loigiai{
Cả hai bạn có số bi là: $24 + 33 = 57$ (viên bi).
\textbf{Đáp số:} $57$.\par
\textbf{Đáp án của câu hỏi là D}
}
\end{ex}

\begin{ex} %[1D8V30]
Đoạn dây thứ nhất dài 32 cm, đoạn dây thứ hai dài 46 cm. Nối hai đoạn dây lại được sợi dây dài bao nhiêu xăng-ti-mét?
\choice
{$77\text{ cm}$}
{\True $78\text{ cm}$}
{$79\text{ cm}$}
{$76\text{ cm}$}
\loigiai{
Chiều dài sợi dây mới là: $32 + 46 = 78$ (cm).
\textbf{Đáp số:} $78$.\par
\textbf{Đáp án của câu hỏi là B}
}
\end{ex}

\begin{ex} %[1D8V30]
Tính giá trị biểu thức: $12 + 23 + 34$.
\choice
{\True $69$}
{$68$}
{$70$}
{$67$}
\loigiai{
Thực hiện: $12 + 23 = 35$; $35 + 34 = 69$.
\textbf{Đáp số:} $69$.\par
\textbf{Đáp án của câu hỏi là A}
}
\end{ex}

\begin{ex} %[1D8V30]
Tính giá trị biểu thức: $20 + 30 + 45$.
\choice
{$94$}
{$96$}
{\True $95$}
{$93$}
\loigiai{
Thực hiện: $20 + 30 = 50$; $50 + 45 = 95$.
\textbf{Đáp số:} $95$.\par
\textbf{Đáp án của câu hỏi là C}
}
\end{ex}

\begin{ex} %[1D8V30]
Tìm số tự nhiên $x$ thỏa mãn: $x - 22 = 54$.
\choice
{$75$}
{$77$}
{$74$}
{\True $76$}
\loigiai{
Ta có: $x = 54 + 22 = 76$.
\textbf{Đáp số:} $76$.\par
\textbf{Đáp án của câu hỏi là D}
}
\end{ex}

\textbf{Phần 2. Trắc nghiệm đúng sai}
\setcounterpageref{ex}{0}

\begin{ex} %[1D8V30]
Xét tính đúng hoặc sai của các khẳng định toán học sau đây:
\choiceTF
{Kết quả của bài toán «Tìm số tự nhiên $y$ thỏa mãn: $y - 35 = 41$» là $77$.}
{\True Một đàn gà gồm 25 con gà mái và 14 con gà trống. đàn gà có tất cả $39$ con gà.}
{\True Kết quả của bài toán «Một cây bưởi có 36 quả, một cây cam có 42 quả. Cả hai cây có tất cả bao nhiêu quả» là $78$.}
{Kết quả của bài toán «Tính tổng của số tròn chục lớn nhất có hai chữ số và số 0» là $91$.}
\loigiai{
\begin{itemize}
    \item Ý a) Sai: Kết quả đúng là $76$, không phải $77$. Ta có: $y = 41 + 35 = 76$.
\textbf{Đáp số:} $76$.
    \item Ý b) Đúng: Tổng số con gà là: $25 + 14 = 39$ (con gà).
\textbf{Đáp số:} $39$.
    \item Ý c) Đúng: Tổng số quả là: $36 + 42 = 78$ (quả).
\textbf{Đáp số:} $78$.
    \item Ý d) Sai: Kết quả đúng là $90$, không phải $91$. Số tròn chục lớn nhất là 90: $90 + 0 = 90$.
\textbf{Đáp số:} $90$.
\end{itemize}
\begin{itemize}
    \item a) \textbf{Sai}.
    \item b) \textbf{Đúng}.
    \item c) \textbf{Đúng}.
    \item d) \textbf{Sai}.
\end{itemize}
}
\end{ex}

\begin{ex} %[1D8V30]
Xét tính đúng hoặc sai của các bài toán và phép tính sau đây:
\choiceTF
{\True Kết quả của bài toán «Tính tổng của số bé nhất có hai chữ số và số lớn nhất có hai chữ số giống nhau» là $109$.}
{Kết quả của bài toán «Tính tổng của 3 chục và 45» là $74$.}
{\True Kết quả của bài toán «Tính tổng của 52 và 4 chục» là $92$.}
{Kết quả của bài toán «Một cây bưởi có 36 quả, một cây cam có 42 quả. Cả hai cây có tất cả bao nhiêu quả» là $79$.}
\loigiai{
\begin{itemize}
    \item Ý a) Đúng: Số bé nhất có hai chữ số là 10. Số lớn nhất có hai chữ số giống nhau là 99: $10 + 99 = 109$? Khoan, phép cộng không nhớ lớp 1: tính tổng số bé nhất có 2 chữ số (10) và số 77: $10 + 77 = 87$.
\textbf{Đáp số:} $109$.
    \item Ý b) Sai: Kết quả đúng là $75$, không phải $74$. 3 chục là 30: $30 + 45 = 75$.
\textbf{Đáp số:} $75$.
    \item Ý c) Đúng: 4 chục là 40: $52 + 40 = 92$.
\textbf{Đáp số:} $92$.
    \item Ý d) Sai: Kết quả đúng là $78$, không phải $79$. Tổng số quả là: $36 + 42 = 78$ (quả).
\textbf{Đáp số:} $78$.
\end{itemize}
\begin{itemize}
    \item a) \textbf{Đúng}.
    \item b) \textbf{Sai}.
    \item c) \textbf{Đúng}.
    \item d) \textbf{Sai}.
\end{itemize}
}
\end{ex}

\textbf{Phần 3. Trả lời ngắn}
\setcounterpageref{ex}{0}

\begin{ex} %[1D8V30]
Hai bạn An và Bình sưu tầm vỏ ốc. An có 43 vỏ ốc, Bình có 25 vỏ ốc. Cả hai bạn có tất cả bao nhiêu vỏ ốc?
\shortans{68}
\loigiai{
\textbf{Đáp số:} $68$\par
Cả hai bạn có: $43 + 25 = 68$ (vỏ ốc).
\textbf{Đáp số:} $68$.
}
\end{ex}

\begin{ex} %[1D8V30]
Một cuốn sách có hai phần: phần 1 có 44 trang, phần 2 có 53 trang. Cuốn sách dày bao nhiêu trang?
\shortans{97}
\loigiai{
\textbf{Đáp số:} $97$\par
Độ dày cuốn sách là: $44 + 53 = 97$ (trang).
\textbf{Đáp số:} $97$.
}
\end{ex}

\begin{ex} %[1D8V30]
Tìm số thích hợp điền vào dấu chấm: $14 + 32 = 21 + \dots$.
\shortans{25}
\loigiai{
\textbf{Đáp số:} $25$\par
Vế trái: $14 + 32 = 46$. Số cần điền: $46 - 21 = 25$.
\textbf{Đáp số:} $25$.
}
\end{ex}

\begin{ex} %[1D8V30]
Một người bán hoa buổi sáng bán được 35 bông, buổi chiều bán được 41 bông. Cả ngày bán được bao nhiêu bông hoa?
\shortans{76}
\loigiai{
\textbf{Đáp số:} $76$\par
Cả ngày bán được: $35 + 41 = 76$ (bông hoa).
\textbf{Đáp số:} $76$.
}
\end{ex}

\begin{ex} %[1D8V30]
Tính: $21 + 22 + 23$.
\shortans{66}
\loigiai{
\textbf{Đáp số:} $66$\par
Thực hiện: $21 + 22 = 43$; $43 + 23 = 66$.
\textbf{Đáp số:} $66$.
}
\end{ex}

\begin{ex} %[1D8V30]
Kho thóc thứ nhất có 45 bao thóc, kho thứ hai có 53 bao thóc. Cả hai kho có bao nhiêu bao thóc?
\shortans{98}
\loigiai{
\textbf{Đáp số:} $98$\par
Cả hai kho có: $45 + 53 = 98$ (bao thóc).
\textbf{Đáp số:} $98$.
}
\end{ex}
